% Shared preamble for all three reports of spec section 17.
%
% Two mechanical rules govern every LaTeX source in this repository, and both
% are enforced in CI rather than by review. They are documented here once,
% because a reader of a chapter will otherwise wonder why the source looks the
% way it does.
%
% THE DASH RULE. Ground rule 13 forbids U+2014 and U+2013 anywhere in the
% project, and forbids the LaTeX ligature forms in .tex sources as well, since
% those typeset to exactly the characters the rule bans.
% scripts/check_dashes.py rejects any run of two or more adjacent hyphens in a
% .tex file. That collides with the one place this project genuinely needs two
% hyphens on the page: a command line flag. The collision is resolved in
% source rather than by an exemption. The macro below inserts an empty group
% between the two hyphens, so the source never carries the adjacent pair and
% the typeset page carries exactly two hyphen glyphs:
%
%     \texttt{\dd{}engine rules}     typesets as the flag a reader would type
%
% Inside a verbatim block the same macro is available because the shell
% environment declares the backslash and braces as command characters.
%
% THE NUMBER RULE. Law 3 forbids a hand typed number anywhere in a report.
% scripts/check_traceability.py scans these sources for anything shaped like a
% measurement and fails unless the line names the committed result file the
% number came from. The chapters of this report therefore contain no digits at
% all: every quoted number is a macro defined in tables/values.tex by
% scripts/make_report_tables.py, beside the path it was read from, and every
% table and figure is generated by that script or by
% scripts/make_report_figures.py. Editing a number by hand is not a policy
% violation that somebody might notice. It is a build failure.

\usepackage[T1]{fontenc}
\usepackage[utf8]{inputenc}
\usepackage{lmodern}
\usepackage[margin=2.5cm]{geometry}
\usepackage{microtype}
\usepackage{booktabs}
\usepackage{longtable}
\usepackage{array}
\usepackage{graphicx}
\usepackage{xcolor}
\usepackage[font=small,labelfont=bf]{caption}
\usepackage{enumitem}
\usepackage{fancyvrb}
\usepackage{parskip}
\usepackage[hidelinks]{hyperref}

% Result paths are printed verbatim and are allowed to break at a slash. The
% url package reads its argument with its own catcodes, so an underscore in a
% run id needs no escape, which matters: an escaped path in the source would no
% longer match the citation pattern the traceability check looks for, and the
% check would silently stop seeing the citation it was given.
%
% Run identifiers are long and carry no slash, so the break list is widened
% inside the command's own style: a run id that cannot break would overflow its
% column rather than wrapping, and the paths are the one thing in these
% documents that a reader may need to type back in.
\DeclareUrlCommand\rpath{%
  \urlstyle{tt}%
  \def\UrlBreaks{\do\/\do\.\do\-\do\_\do\T\do\Z\do\p\do\v\do\b\do\a}%
  \def\UrlBigBreaks{\do\/}%
}

% An identifier printed verbatim in a narrow table column: a metric name, a
% slice name, a taxonomy label. Same mechanism as \rpath, with breaks allowed
% at the separators identifiers actually use, so a long label wraps inside its
% column instead of running off the page.
\DeclareUrlCommand\ident{%
  \urlstyle{tt}%
  \def\UrlBreaks{\do\_\do\-\do\/\do\:\do\.}%
  \def\UrlBigBreaks{}%
}

% The two hyphen macro described at the top of this file.
\newcommand{\dd}{-{}-}

% A command line flag, written as the reader would type it.
\newcommand{\flag}[1]{\texttt{\dd{}#1}}

% Console transcripts. commandchars lets \dd{} through so that a flag can be
% shown inside an otherwise verbatim block.
\DefineVerbatimEnvironment{shell}{Verbatim}{%
  commandchars=\\\{\},%
  fontsize=\small,%
  frame=single,%
  framesep=3mm,%
  rulecolor=\color{black!30}%
}

% Markup samples, with no command characters at all, so angle brackets, quotes
% and braces are exactly what they look like.
\DefineVerbatimEnvironment{markup}{Verbatim}{%
  fontsize=\small,%
  frame=single,%
  framesep=3mm,%
  rulecolor=\color{black!30}%
}

% A finding code, a taxonomy label, or any other identifier that has one exact
% spelling and must not be hyphenated across a line break.
\newcommand{\code}[1]{\texttt{#1}}

% The standing caveat, printed wherever a reader could otherwise mistake a
% measurement on this corpus for a claim about the open web.
\newcommand{\syntheticnote}{%
  \emph{Measured on a seeded synthetic corpus generated by this repository.
  Behaviour on real world pages is unmeasured and is stated as unmeasured.}%
}

\setlist{itemsep=2pt,parsep=0pt,topsep=4pt}
\captionsetup[table]{skip=4pt}
\captionsetup[figure]{skip=6pt}
